\documentclass[11pt,a4paper]{article}
\usepackage[UTF8]{ctex}
\usepackage{geometry}
\geometry{left=2.5cm,right=2.5cm,top=2.5cm,bottom=2.5cm}
\usepackage{amsmath,amssymb,amsthm}
\usepackage{hyperref}
\usepackage{booktabs}
\usepackage{enumitem}
\usepackage{xltabular}
\hypersetup{colorlinks=true,linkcolor=blue,urlcolor=blue}

\newtheorem{definition}{定义}[section]
\newtheorem{proposition}{命题}[section]
\newtheorem{remark}{注记}[section]

\title{基于SMPC的协同逼移与对抗\\
——密码学武装下的攻击方联盟与防御方反制}
\author{李龙胜}
\date{\today}

\begin{document}

\maketitle

\begin{center}
\textit{
当攻击方也掌握了安全多方计算，会发生什么？\\
他们可以在不彼此暴露资源的前提下，\\
联合侦察、虚拟化IP池、协同逼移。\\
防御方的信息论下界被绕过。\\
但SMPC也是双刃剑——它依赖的公开参与和隐私保护，\\
恰好是防御方投毒和瓦解联盟的入口。\\
密码学没有解决信任问题，它只是把信任问题转移到了新的层面。
}
\end{center}

\vspace{5mm}

\begin{abstract}
此前关于多攻击方协同的讨论聚焦于物理层面的资源协作
（信息共享、分工、轮班）。
然而，攻击方可以引入安全多方计算（SMPC），
在不彼此暴露私有信息的前提下联合执行最优逼移攻击。
本文建立基于SMPC的攻击方联盟协议，
推导联盟如何绕过防御方的信息论下界，
虚拟化IP资源，以及匿名化攻击任务市场。
在防御方一侧，提出SMPC对抗SMPC的防御策略——
构建防御方SMPC联盟进行联合威胁分析，
以及利用投毒和女巫攻击从内部瓦解攻击方联盟。
\end{abstract}

\tableofcontents

\section{攻击方SMPC协议的基本设定}

设存在 \(M\) 个攻击方 \(A_1, A_2, \dots, A_M\)，
每个攻击方 \(A_m\) 拥有私有输入：
\begin{itemize}
    \item 资源预算 \(C_m\)（算力、IP池）。
    \item 对防御方参数的局部信念 \(I_{\text{def}}^{(m)}\)（基于自身历史探测数据）。
    \item 攻击意图 \(G_m\)（目标资产、期望的攻击效果）。
\end{itemize}

攻击方联盟的目标是协同执行一次最优逼移攻击，
但不彼此泄露各自的私有输入。
他们通过一个 SMPC 协议，安全地计算联合函数：
\[
(i_{\text{b}}^*, i_{\text{a}}^*, \{n_m^*\}, \{t_m^*\}) = F(\{C_m\}, \{I_{\text{def}}^{(m)}\}, \{G_m\})
\]
其中：
\begin{itemize}
    \item \(i_{\text{b}}^*\) 是最优逼移方向。
    \item \(i_{\text{a}}^*\) 是最优真实攻击方向。
    \item \(\{n_m^*\}\) 是分配给各攻击方的最优逼移量。
    \item \(\{t_m^*\}\) 是最优攻击时机序列。
\end{itemize}

协议结束后，每个攻击方仅知道自己的任务切片，
而对联盟的整体攻击计划和盟友的私有信息一无所知。
信任不再依赖攻击方之间的脆弱承诺，
而是锚定在SMPC的密码学保证上。

\section{攻击方联盟的优势升级}

SMPC将攻击方联盟的优势从物理协同升级为信息论协同。

\subsection{突破防御方信息论下界}

此前我们证明：防御方可以通过控制单周期最大探测限制 \(M_{\text{period}}\)
来使攻击方无法获得足够样本。
SMPC 使攻击方得以绕过这一限制。

假设每个攻击方对某类告警 \(i\) 的探测结果为
\(X_m \sim \text{Binomial}(n_m, P_{\text{obs}})\)。
攻击方联盟通过 SMPC 安全地聚合出总观测样本：
\[
X = \sum_{m=1}^{M} X_m, \quad n = \sum_{m=1}^{M} n_m
\]
每个攻击方仅知道自己的 \(X_m\) 和 \(n_m\)，
但联盟共同获得了联合估计 \(\hat{P}_{\text{obs}} = X/n\)。

联合有效样本量 \(n\) 可以轻松超越 \(M_{\text{period}}\)。
防御方为单攻击方设计的信息论下界被绕过。
代价是攻击方联盟需要为执行 SMPC 协议支付额外的计算和通信成本。

\subsection{虚拟化IP资源}

通过 SMPC，攻击方可以在不暴露 IP 所有权的前提下，
将闲置 IP 借给其他攻击方用于发动逼移。
攻击方 A 生成一组加密的代理令牌，
只有指定的攻击方 B 能在有效期内使用这些令牌。
A 全程不知道 B 的攻击目标，
B 也完全不知道令牌的真正来源。

整个IP资源在联盟内部形成虚拟的、可动态调度的暗云。
IP池不再是分散的、静态的，
而是统一的、按需分配的虚拟资源。

\subsection{匿名攻击任务市场}

在暗云IP池之上，可以构建匿名的攻击任务市场。
任一攻击方可以匿名发布任务：
攻击类型、时间窗口、目标范围、赏金（加密货币）。
闲置的攻击方认领任务并执行，
通过 SMPC 提交任务完成的零知识证明来领取赏金。

赏金由智能合约托管，
既不知道发布者是谁也不知道执行者是谁。
这是暗网的终极形态——
不是买卖数据的市场，而是买卖攻击力本身的市场。

\section{防御方对策}

面对这一威胁，防御方同样需要引入 SMPC，
并采取新的博弈策略。

\subsection{构建防御方SMPC联盟}

多家企业通过 SMPC 共同计算联合威胁模型，
识别当前攻击方联盟的逼移模式。
各家企业的台账参数（\(A_i^{\min}, \alpha_i\)）保持加密，
仅共享为了计算最优防御策略所必需的最小信息。

当某企业被攻击时，
它可以将攻击流量特征在加密状态下广播给联盟。
其他成员在不暴露自身网络结构的前提下，
在自己的日志中匹配并协助溯源。

\subsection{从内部瓦解攻击方联盟}

防御方无需破解 SMPC 的密码学保障，
而是攻击其博弈论脆弱性。

\textbf{投毒污染}：
防御方将虚假的内部破坏者渗透进攻击方联盟，
提供精心设计的污染样本。
这些样本通过 SMPC 的隐私保护机制混入联盟的联合侦察数据中，
在安全聚合过程中无法被识别出来。
防御方可以诱导攻击方高估某个方向的防御薄弱度，
从而将逼移攻击引向预设的陷阱。

\textbf{女巫淹没}：
防御方创建大量匿名身份加入攻击方联盟。
在需要投票或协商的决策环节，
用数量优势控制决策结果。
攻击方的SMPC协议若依赖多数诚实假设，
则女巫攻击将构成致命威胁。

\section{讽刺性闭环}

攻击方通过 SMPC 解决了内部信任问题，
信任从脆弱的承诺转移到密码学上。
但防御方的投毒和女巫攻击，
恰好利用了SMPC的匿名性和开放性。

信任问题没有消失，
它只是从攻击方之间的承诺层，
转移到了攻击方对联盟成员身份的验证层。
旧的脆弱性被密码学缝合了，
新的脆弱性在更深的层次上裂开。

\section{诚实标注}

\begin{center}
\begin{xltabular}{\textwidth}{c|X|c}
\toprule
\textbf{命题} & \textbf{状态} & \textbf{层级} \\
\midrule
攻击方SMPC协议设计 &
理想函数框架已给出，具体实现依赖成熟密码学工具 &
II \\
联合侦察突破信息论下界 &
统计聚合逻辑已验证，SMPC工程实现可行 &
II \\
虚拟IP池与暗云 &
基于现有隐藏服务技术的推演，技术可行性存在 &
III \\
匿名攻击任务市场 &
概念完整，需要匿名支付系统与SMPC的深度集成 &
III \\
防御方SMPC联盟 &
设计已给出，大规模部署的非技术壁垒极高 &
III \\
投毒与女巫反制 &
逻辑上可行，实战效果待验证 &
III \\
\bottomrule
\end{xltabular}
\end{center}

\section{结语}

攻击方引入 SMPC 进行协同逼移，
标志着密码学正式成为网络攻击的基础设施。
防御方的信息论下界被绕过，
IP池被虚拟化，攻击力本身被市场化。
而防御方的反制路径同样清晰——
用 SMPC 对抗 SMPC，
从内部瓦解攻击方联盟。
攻防双方在密码学的军备竞赛中共同升级。

\vspace{5mm}
\noindent\textbf{文档信息}：
\begin{itemize}
    \item 作者：李龙胜
    \item 版本：第一版（SMPC协同逼移归档）
    \item 许可：CC BY 4.0
    \item 前置依赖：多攻击方协同与竞争、密码学融合蓝图、安全参数自适应调节
\end{itemize}

\end{document}